\documentclass[10pt,a4paper]{amsart}
\usepackage[margin=19mm]{geometry}
\usepackage{amsmath,amssymb,mathtools}
\usepackage[scr=rsfs,cal=euler]{mathalfa}
\usepackage{xfrac}
\usepackage[unicode,hidelinks,bookmarks=false]{hyperref}
\newcommand{\ie}{i.e.\ }
\newcommand{\cf}{cf.\ }
\newcommand{\resp}{respectively\ }
\newcommand{\Subsection}{\subsection}
\providecommand{\nobreakdash}{\nobreak\hbox{-}}
\setlength{\emergencystretch}{2em}
\multlinegap0pt
\usepackage{bm}
\usepackage{bbm}
\usepackage{dsfont}
\usepackage{xspace}
\usepackage{cancel}
\usepackage{mathtools}
\usepackage{amsthm}
\makeatletter
\@ifundefined{theorem}{\newtheorem{theorem}{Theorem}}{}
\@ifundefined{lemma}{\newtheorem{lemma}[theorem]{Lemma}}{}
\@ifundefined{proposition}{\newtheorem{proposition}[theorem]{Proposition}}{}
\makeatother
\newcommand{\HS}{\mathrm{HS}}
\newcommand{\NN}{\mathbb{N}}
\title[Spectral Structure of the Mixed Hessian of the Dispersionles]{Spectral Structure of the Mixed Hessian of the Dispersionless Toda $\tau$-Function}
\author{Oleg Alekseev}
\date{Source: arXiv:2603.23424v2 (CC BY 4.0)}
\begin{document}
\maketitle

\noindent\emph{Source and licence.} Spectral Structure of the Mixed Hessian of the Dispersionless Toda $\tau$-Function. Version arXiv:2603.23424v2 (CC BY 4.0), distributed under the Creative Commons Attribution 4.0 International licence, as recorded in the arXiv licence metadata for this exact version and shown on the abstract page https://arxiv.org/abs/2603.23424v2. Full rights notice: licenses/arxiv-2603.23424-CC-BY-4.0.txt.\par\smallskip

\noindent\textit{Editorial scope.} Counted unit: the Lemma whose source label is \texttt{lem:entrywise-log} in the article (source0.tex lines 3639--3664, source label \texttt{lem:entrywise-log}), delivered together with the article's own complete original proof of it (source0.tex lines 3666--3879, one \texttt{proof} environment of 214 lines). No mathematics is written, repaired or re-derived: statement and proof are the article's own text line for line. Required context retained uncounted: eq:app-Raney-asymp (lines 3340--3351); prop:raney-uniform (lines 3392--3398); lem:raney-uniform-expansion (lines 3456--3467); lem:Gram-entry (lines 3534--3550); eq:app-d-def (lines 3586--3589); lem:raney-gap (lines 3593--3601). Every definition, display and label of those lines is retained unchanged. Omitted from this package and not redistributed: 1--3339, 3352--3391, 3399--3455, 3468--3533, 3551--3585, 3590--3592, 3602--3638, 3665--3665, 3880--4136. Nothing inside a retained excerpt is omitted or altered: every retained body is delivered exactly as the article's lines are written, so no omission receipt of any kind accompanies this package. Citations: the retained excerpts carry no citation command at all, so no bibliography entry of the article is transcribed and no reference is redistributed. Reference adaptation: none was needed and none was made; every reference inside the delivered unit resolves inside it, so no unresolved reference or citation is produced. Printed slips kept literal: no printed word, symbol, hypothesis or number of the counted statement or of its proof was altered. Layout and class: the article's own class is \texttt{amsart} with packages including fontenc, microtype, amsmath, amssymb, amsfonts, amsthm, mathtools, mathrsfs, bm, placeins, geometry, graphicx, enumitem, natbib; the wrapper is the lane's pinned \texttt{amsart} editorial wrapper at 10pt on A4, which restates the theorem-like environments the retained text needs and adds the article's own macro definitions that the retained text needs (\texttt{\textbackslash HS}, \texttt{\textbackslash NN}), with extra wrapper packages (bm, bbm, dsfont, xspace, cancel, mathtools). The delivered numbering is the unit's own. Assets: no retained span contains a figure environment, a graphics-inclusion command, a PDF-page inclusion, a table or a TikZ picture; no image file is copied into this package, the package declares no asset dependency and no image file is redistributed with it. Licence: version arXiv:2603.23424v2 of this article is distributed under the Creative Commons Attribution 4.0 International licence, as recorded in the arXiv licence metadata for this exact version and shown on the abstract page https://arxiv.org/abs/2603.23424v2. A full rights notice is delivered at licenses/arxiv-2603.23424-CC-BY-4.0.txt. Characters: every retained excerpt is pure ASCII, so the delivered engine, XeLaTeX with its default Latin Modern text font, reports no missing character.
	\begin{equation}\label{eq:app-Raney-asymp}
		R_{s,p}(m)
		=
		A_{s,p}\,\zeta_c^{-m}\,m^{-3/2}\bigl(1+O(m^{-1})\bigr),
		\qquad
		A_{s,p}
		=
		\frac{p}{\sqrt{2\pi}}\,
		\frac{s^{\,p-\frac12}}{(s-1)^{\,p+\frac12}}
		=
		\frac{p\,M^{p}}{\sqrt{2\pi s(s-1)}}.
	\end{equation}

\begin{proposition}[Uniform upper bound]\label{prop:raney-uniform}
	There exists a constant $C_s>0$ such that for all integers $p\ge1$ and $m\ge1$,
	\begin{equation}\label{eq:raney-uniform}
		R_{s,p}(m)\le C_s\,p\,M^{p}\,\zeta_c^{-m}\,m^{-3/2}.
	\end{equation}
	Moreover, the same bound holds with $m^{-3/2}$ replaced by $(1+m)^{-3/2}$ for all $m\ge0$.
\end{proposition}

\begin{lemma}[Uniform one-term expansion in the tail region]\label{lem:raney-uniform-expansion}
	Fix \(s\ge2\) and \(0<\theta<1\).  There exists a constant \(C_{s,\theta}>0\)
	such that for all integers \(p\ge1\) and \(m\ge1\) satisfying \(p\le \theta m\),
	\begin{equation}\label{eq:raney-uniform-expansion}
		R_{s,p}(m)
		=
		A_{s,p}\,\zeta_c^{-m}\,m^{-3/2}\Bigl(1+\varepsilon_{s,p}(m)\Bigr),
		\qquad
		|\varepsilon_{s,p}(m)|\le \frac{C_{s,\theta}}{m},
	\end{equation}
	where $A_{s,p}$ is as in \eqref{eq:app-Raney-asymp}.
\end{lemma}

\begin{lemma}[Entries of the Gram block]\label{lem:Gram-entry}
	For $j_1\le j_2$ (and $\Delta:=j_2-j_1\ge0$) one has
	\begin{equation}\label{eq:app-Gram-entry}
		G^{(q)}_{j_1j_2}(\zeta)
		=
		\sum_{m=0}^{\infty}
		\frac{(p_{j_2}+sm)^{2}}{\sqrt{p_{j_1}p_{j_2}}}\,
		R_{s,p_{j_1}}(m+\Delta)\,R_{s,p_{j_2}}(m)\,\zeta^{\,2m+\Delta}.
	\end{equation}
	Consequently,
	\begin{equation}\label{eq:app-Gtilde-entry-clean}
		\widetilde G^{(q)}_{j_1j_2}(\zeta)
		=
		\frac{1}{w_{j_1}w_{j_2}}\,G^{(q)}_{j_1j_2}(\zeta),
		\qquad j_1,j_2\in\NN_0.
	\end{equation}
\end{lemma}

\begin{equation}\label{eq:app-d-def}
	d^{(q)}_j:=\frac{s\,A_{s,p_j}}{\sqrt{p_j}},\qquad \widetilde d^{(q)}_j:=\frac{d^{(q)}_j}{w_j},
	\qquad j\in\NN_0.
\end{equation}

\begin{lemma}[Uniform exponential gap away from the critical slope]\label{lem:raney-gap}
	Fix $s\ge2$ and $\theta>0$.  There exist constants $C_{s,\theta},c_{s,\theta}>0$
	such that for all integers $p,m\ge1$ satisfying $p\ge\theta m$,
	\[
		R_{s,p}(m)
		\le
		C_{s,\theta}\,p\,M^{p}\,\zeta_c^{-m}\,m^{-3/2}\,e^{-c_{s,\theta}p}.
	\]
\end{lemma}

\begin{lemma}[Logarithmic tail with Hilbert--Schmidt remainder]\label{lem:entrywise-log}
	Let $\eta=\zeta/\zeta_c\in(0,1)$ and $L(\zeta)=\log\frac{1}{1 - \eta^2}$.
	Define the rank--one Toeplitz kernel
	\[
		(K_\eta)_{j_1j_2}:=\eta^{|j_1-j_2|}\,\widetilde d^{(q)}_{j_1}\,\widetilde d^{(q)}_{j_2},
		\qquad j_1,j_2\in\NN_0.
	\]
	Then, for every $0<\zeta<\zeta_c$, the operator
	\[
	R^{(q)}(\zeta):=\widetilde G^{(q)}(\zeta)-L(\zeta)\,K_\eta
	\]
	is Hilbert--Schmidt on $\ell^2(\NN_0)$, with $\sup_{\zeta<\zeta_c}\|R^{(q)}(\zeta)\|_{\mathrm{HS}}<\infty$.
 Moreover, as $\zeta\uparrow\zeta_c$, $R^{(q)}(\zeta)$ converges in Hilbert--Schmidt norm
 to a limit Hilbert--Schmidt operator $R_*^{(q)}$, hence also in operator norm.

	In particular, for each fixed $(j_1,j_2)$ one has the entrywise expansion
	\begin{equation}\label{eq:app-entrywise-log}
		\widetilde G^{(q)}_{j_1j_2}(\zeta)
		=
		\widetilde d^{(q)}_{j_1}\,\widetilde d^{(q)}_{j_2}\,\eta^{|j_1-j_2|}\,L(\zeta)
		+(R_*^{(q)})_{j_1j_2}
		+o(1),
		\qquad \zeta\uparrow\zeta_c,
	\end{equation}
	and the error is $o(1)$ in $\ell^2$--matrix (Hilbert--Schmidt) sense.
\end{lemma}

\begin{proof}
	By symmetry of the Gram kernel, it is enough to treat the case \(j_1\le j_2\). Write \(\Delta:=j_2-j_1\ge0\).  By
	Lemma~\ref{lem:Gram-entry},
	\[
		\widetilde G^{(q)}_{j_1j_2}(\zeta)=\sum_{m\ge0}U_m(j_1,j_2;\zeta),
	\]
	where
	\[
		U_{m}(j_1,j_2;\zeta)
		:=
		\frac{1}{w_{j_1}w_{j_2}}\,
		\frac{(p_{j_2}+sm)^{2}}{\sqrt{p_{j_1}p_{j_2}}}\,
		R_{s,p_{j_1}}(m+\Delta)\,R_{s,p_{j_2}}(m)\,\zeta^{\,2m+\Delta}.
	\]
	Since \(p_{j_2}\ge p_{j_1}\) and \(p_{j_2}\ge \Delta\), the cutoff
	\[
		M_{j_1j_2}:=\max\{2p_{j_1},\,2p_{j_2},\,2\Delta,\,1\}
	\]
	reduces here to \(M_{j_1j_2}=2p_{j_2}\).  We split the sum into the initial
	range \(0\le m<2p_{j_2}\) and the tail \(m\ge2p_{j_2}\).


	\smallskip
	\noindent\emph{Tail region.}
	If \(m\ge2p_{j_2}\), then \(p_{j_2}\le m/2\) and
	\(p_{j_1}\le (m+\Delta)/2\), so Lemma~\ref{lem:raney-uniform-expansion}
	applies to both Raney factors (with \(\theta=\frac12\)). More explicitly,
	\[
		R_{s,p_{j_1}}(m+\Delta)
		=
		A_{s,p_{j_1}}\zeta_c^{-m-\Delta}(m+\Delta)^{-3/2}
		\Bigl(1+O((m+1)^{-1})\Bigr),
	\]
	and
	\[
		R_{s,p_{j_2}}(m)
		=
		A_{s,p_{j_2}}\zeta_c^{-m}m^{-3/2}
		\Bigl(1+O((m+1)^{-1})\Bigr),
	\]
	uniformly for \(m\ge2p_{j_2}\). Since \(j_1\le j_2\), the quantities
	\(p_{j_1},p_{j_2},\Delta\) are fixed along the tail sum, while
	\[
		(p_{j_2}+sm)^2=s^2m^2\Bigl(1+O((m+1)^{-1})\Bigr)
	\]
	and
	\[
		(m+\Delta)^{-3/2}m^{-3/2}
		=
		m^{-3}\Bigl(1+O((m+1)^{-1})\Bigr).
	\]
	Hence
	\[
		\frac{(p_{j_2}+sm)^2}{\sqrt{p_{j_1}p_{j_2}}}\,
		R_{s,p_{j_1}}(m+\Delta)\,R_{s,p_{j_2}}(m)\,\zeta^{2m+\Delta}
		=
		\frac{s^2A_{s,p_{j_1}}A_{s,p_{j_2}}}{\sqrt{p_{j_1}p_{j_2}}}\,
		\eta^\Delta\frac{\eta^{2m}}{m}
		\Bigl(1+O((m+1)^{-1})\Bigr).
	\]
	Using \eqref{eq:app-d-def} and \(m^{-1}=(m+1)^{-1}+O((m+1)^{-2})\), we
	obtain
	\begin{equation}\label{eq:Sm-decomp}
		U_m(j_1,j_2;\zeta)
		=
		\widetilde d^{(q)}_{j_1}\widetilde d^{(q)}_{j_2}\,\eta^{\Delta}\,
		\frac{\eta^{2m}}{m+1}
		+E_m^{\mathrm{tail}}(j_1,j_2;\zeta),
	\end{equation}
	with
	\begin{equation}\label{eq:Sm-error}
		|E_m^{\mathrm{tail}}(j_1,j_2;\zeta)|
		\le
		C\,\widetilde d^{(q)}_{j_1}\widetilde d^{(q)}_{j_2}\,\eta^{\Delta}\,
		\frac{\eta^{2m}}{(m+1)^2},
		\qquad m\ge 2p_{j_2},
	\end{equation}
	where \(C\) depends only on \(s\) and \(\beta\).  Summing over \(m\ge2p_{j_2}\)
	shows that
	\[
		E^{\mathrm{tail}}_{j_1j_2}(\zeta)
		:=\sum_{m\ge2p_{j_2}}E_m^{\mathrm{tail}}(j_1,j_2;\zeta)
	\]
	defines a uniformly Hilbert--Schmidt kernel, since
	\[
		|E^{\mathrm{tail}}_{j_1j_2}(\zeta)|
		\le
		C\,\widetilde d^{(q)}_{j_1}\widetilde d^{(q)}_{j_2}
		\sum_{m\ge0}\frac{1}{(m+1)^2}
		\le
		C'\,\widetilde d^{(q)}_{j_1}\widetilde d^{(q)}_{j_2},
	\]
	and \(\widetilde{\bm d}^{(q)}\in\ell^2(\NN_0)\).


	\smallskip
	\noindent\emph{Initial range.}
	Write
	\[
		E^{\mathrm{init}}_{j_1j_2}(\zeta):=\sum_{0\le m<2p_{j_2}}U_m(j_1,j_2;\zeta).
	\]
	The term \(m=0\) is estimated by Proposition~\ref{prop:raney-uniform}:
	\[
		|U_0(j_1,j_2;\zeta)|
		\le
		C\,M^{-p_{j_2}}\,p_{j_1}^{-1-\beta}p_{j_2}^{1-\beta}(1+\Delta)^{-3/2}.
	\]
	For \(1\le m<2p_{j_2}\), Proposition~\ref{prop:raney-uniform} applies to
	\(R_{s,p_{j_1}}(m+\Delta)\), while Lemma~\ref{lem:raney-gap} with
	\(\theta=\frac12\) applies to \(R_{s,p_{j_2}}(m)\) because
	\(p_{j_2}\ge m/2\).  Using \(\zeta^{2m+\Delta}\le \zeta_c^{2m+\Delta}\), one gets
	\[
		|U_m(j_1,j_2;\zeta)|
		\le
		C\,e^{-c p_{j_2}}\,
		p_{j_1}^{-1-\beta}p_{j_2}^{-1-\beta}\,
		\frac{(p_{j_2}+sm)^2}{(m+1)^{3/2}(m+\Delta+1)^{3/2}}.
	\]
	Since \(m<2p_{j_2}\), we have \(p_{j_2}+sm\le (1+2s)p_{j_2}\), and therefore
	\[
		\sum_{m=1}^{2p_{j_2}-1}|U_m(j_1,j_2;\zeta)|
		\le
		C\,e^{-c p_{j_2}}\,
		p_{j_1}^{-1-\beta}p_{j_2}^{1-\beta}
		\sum_{m\ge1}\frac{1}{(m+1)^3}
		\le
		C'\,e^{-c p_{j_2}}\,p_{j_1}^{-1-\beta}p_{j_2}^{1-\beta}.
	\]
	Hence
	\[
		|E^{\mathrm{init}}_{j_1j_2}(\zeta)|
		\le
		C''\,e^{-c' p_{j_2}}\,p_{j_1}^{-1-\beta}p_{j_2}^{1-\beta},
	\]
	for constants \(C'',c'>0\) independent of \(\zeta\).  This bound is square
	summable in \((j_1,j_2)\), so \(E^{\mathrm{init}}(\zeta)\) is uniformly
	Hilbert--Schmidt.  Moreover, for each fixed \((j_1,j_2)\) the sum is finite and
	therefore converges termwise as \(\zeta\uparrow\zeta_c\).

	Summing \eqref{eq:Sm-decomp} over \(m\ge2p_{j_2}\) gives
	\[
		\sum_{m\ge2p_{j_2}}
		U_m(j_1,j_2;\zeta)
		=
		\widetilde d^{(q)}_{j_1}\widetilde d^{(q)}_{j_2}\,\eta^\Delta
		\sum_{m\ge2p_{j_2}}\frac{\eta^{2m}}{m+1}
		+
		E^{\mathrm{tail}}_{j_1j_2}(\zeta).
	\]
	Reinstating the omitted harmonic terms produces the correction
	\[
		E^{\mathrm{harm}}_{j_1j_2}(\zeta)
		:=
		\widetilde d^{(q)}_{j_1}\widetilde d^{(q)}_{j_2}\,\eta^\Delta
		\sum_{m=0}^{2p_{j_2}-1}\frac{\eta^{2m}}{m+1}.
	\]
	Since
	\[
		\sum_{m=0}^{2p_{j_2}-1}\frac{\eta^{2m}}{m+1}
		\le
		1+\log(1+2p_{j_2}),
	\]
	we obtain
	\[
		|E^{\mathrm{harm}}_{j_1j_2}(\zeta)|
		\le
		C\,\widetilde d^{(q)}_{j_1}\widetilde d^{(q)}_{j_2}\bigl(1+\log(1+p_{j_2})\bigr).
	\]
	Because \(\widetilde d^{(q)}_j\asymp p_j^{-1-\beta}\), the kernel on the
	right-hand side is Hilbert--Schmidt:
	\[
		\sum_{j_1,j_2\ge0}
		\widetilde d^{(q)}_{j_1}{}^2\widetilde d^{(q)}_{j_2}{}^2
		\bigl(1+\log(1+p_{j_2})\bigr)^2
		<\infty.
	\]

	Thus
	\[
		\widetilde G^{(q)}_{j_1j_2}(\zeta)
		=
		\widetilde d^{(q)}_{j_1}\widetilde d^{(q)}_{j_2}\,\eta^\Delta
		\sum_{m\ge0}\frac{\eta^{2m}}{m+1}
		+
		E^{\mathrm{init}}_{j_1j_2}(\zeta)
		+
		E^{\mathrm{tail}}_{j_1j_2}(\zeta)
		+
		E^{\mathrm{harm}}_{j_1j_2}(\zeta).
	\]
	Since \(\sum_{m\ge0}\frac{\eta^{2m}}{m+1}=\eta^{-2}L(\zeta)\), we write
	\[
		\eta^{-2}L(\zeta)=L(\zeta)+\bigl(\eta^{-2}-1\bigr)L(\zeta).
	\]
	The coefficient \(\bigl(\eta^{-2}-1\bigr)L(\zeta)\) stays bounded as
\(\eta\uparrow1\), and \(K_\eta\) is uniformly Hilbert--Schmidt. Therefore the
extra term $\bigl(\eta^{-2}-1\bigr)L(\zeta)\,K_\eta$
is uniformly Hilbert--Schmidt as well. Collecting all Hilbert--Schmidt pieces
into \(R^{(q)}(\zeta)\) yields
\[
    \widetilde G^{(q)}(\zeta)=L(\zeta)\,K_\eta+R^{(q)}(\zeta),
\]
with \(\sup_{\zeta<\zeta_c}\|R^{(q)}(\zeta)\|_{\HS}<\infty\).

 Finally, each constituent of \(R^{(q)}(\zeta)\) has a pointwise
	limit as \(\zeta\uparrow\zeta_c\), and the dominating Hilbert--Schmidt kernels
	displayed above are independent of \(\zeta\).  Dominated convergence therefore
	yields
	\[
	 R^{(q)}(\zeta)\to R_*^{(q)}
		\qquad\text{in Hilbert--Schmidt norm as }\zeta\uparrow\zeta_c,
	\]
	hence also in operator norm.
\end{proof}

\end{document}
